\section{Comparadores, bits y aleatoriedad}

\subsection{Orden personalizado}

El comparador debe ser estricto: nunca puede decir que $a<a$. No usar una
tolerancia EPS como comparador directo de \texttt{set} o \texttt{sort}.

\begin{lstlisting} [style=cpp]
struct BySecondThenFirst {
    bool operator()(const pair<int,int> &a, const pair<int,int> &b) const {
        return tie(a.second, a.first) < tie(b.second, b.first);
    }
};

void comparator_example(vector<pair<int,int>> &items) {
    sort(items.begin(), items.end(), BySecondThenFirst{});
    sort(items.begin(), items.end(), [](const auto &a, const auto &b) {
        if (a.first != b.first) return a.first > b.first;
        return a.second < b.second;
    });
}
\end{lstlisting}

\subsection{Operaciones de bits}

\begin{center}
\begin{tabular}{|c|c|c|}
\hline
Consulta & \texttt{oper} & \texttt{NEUT} \\
\hline
XOR & \texttt{\^} & \texttt{0} \\
OR & \texttt{|} & \texttt{0} \\
AND & \texttt{\&} & \texttt{(\textasciitilde 0)} \\
\hline
\end{tabular}
\end{center}

\begin{lstlisting}[style=cpp]
int count_bits(unsigned long long value) {
    return __builtin_popcountll(value);
} //cuenta cuantos bits estan en 1

int least_significant_bit_index(unsigned long long value) {
    assert(value != 0);
    return __builtin_ctzll(value);
} // devuelve cuantos ceros hay al final, o sea el indice del bit 1 mas bajo

int most_significant_bit_index(unsigned long long value) {
    assert(value != 0);
    return 63 - __builtin_clzll(value);
} // obtiene el indice del bit 1 mas alto

unsigned long long lowest_set_bit(unsigned long long value) {
    return value & (~value + 1); // tambien value & -value para unsigned
} // devuelve solo el bit 1 mas bajo

void bitset_example() {
    bits.set(5);    // pone el bit 5 en 1
    bits.flip(7);   // invierte el bit 7
    bits.reset(5);  // pone el bit 5 en 0
    bits.test(7);   // consulta si el bit 7 esta en 1
    bits.count();   // cuenta bits en 1
    (void)seventh; (void)active;
}
\end{lstlisting}

\warning{\texttt{\_\_builtin\_clz(0)} y \texttt{\_\_builtin\_ctz(0)} son
indefinidos. Para \texttt{long long}, usar las versiones terminadas en
\texttt{ll}.}

\subsection{Random y shuffle}

\begin{lstlisting}[style=cpp]
mt19937_64& random_engine() {
    static mt19937_64 engine(chrono::steady_clock::now()
                             .time_since_epoch().count());
    return engine;
}

ll random_integer(ll left, ll right) {
    return uniform_int_distribution<ll>(left, right)(random_engine());
}

void shuffle_values(vector<int> &values) {
    shuffle(values.begin(), values.end(), random_engine());
}
\end{lstlisting}

\subsection{Errores frecuentes}

\begin{itemize}
  \item \texttt{map[key]} y \texttt{unordered\_map[key]} insertan la clave si
  no existe; usar \texttt{find} para consultar sin modificar.
  \item No desreferenciar \texttt{end()}, ni llamar \texttt{front/back/top}
  sobre un contenedor vac\'io.
  \item \texttt{erase(value)} en \texttt{multiset} borra todas las copias;
  \texttt{erase(find(value))} borra una.
  \item Un comparador de \texttt{sort/set/map} debe cumplir orden estricto.
  \item Evitar \texttt{endl} salvo que se necesite forzar el flush; usar
  \texttt{'\textbackslash n'}.
  \item En expresiones modulares, normalizar negativos con
  $((x\bmod m)+m)\bmod m$.
\end{itemize}
